% ======================================================================
% Appendix A: Data construction (draft, written from the replication
% package; every number below is taken from replication/data/moments,
% replication/output or the package logs).
% ======================================================================
\section{Data Construction}\label{app:data}


This appendix documents the samples, definitions, and calculations behind
the facts in Section~\ref{sec:data} and the estimation targets in
Section~\ref{sec:estimation}, in the order in which the main text uses
them. Subsection~\ref{app:data-ehs} describes the English Housing Survey
files and definitions, Subsection~\ref{app:data-regression} the
age-gradient regressions, Subsection~\ref{app:data-epc} the EPC register,
Subsections~\ref{app:data-was} and~\ref{app:data-ukhls} the Wealth and
Assets Survey and Understanding Society, Subsection~\ref{app:data-other}
mortality and energy bills, and Subsection~\ref{app:data-targets} the
samples behind each estimation target.

% ----------------------------------------------------------------------
\subsection{English Housing Survey: files and definitions.}\label{app:data-ehs}

\paragraph{Files.} The EHS is released in two forms. A housing-stock release
pools two consecutive years of fieldwork and contains the households whose
dwelling received a physical inspection, together with their interview
records and a weight for the pooled sample. A fieldwork release contains the
full interview sample of a single year, without physical data.
Table~\ref{tab:ehs-releases} lists the releases we use. Statistics that need
the physical survey (the EPC band, the SAP rating, and dwelling age) use
housing-stock releases only.

\begin{table}[htbp]
  \centering
  \caption{English Housing Survey releases}\label{tab:ehs-releases}
  \begin{threeparttable}
    \small
    \setlength{\tabcolsep}{4pt}
    \begin{tabular}{@{}llllccc@{}}
      \toprule
      & & & & \multicolumn{3}{c}{Used for} \\
      \cmidrule(l){5-7}
      Release & File & UKDA SN & Fieldwork years & Homeownership & A--C share & Regression \\
      \midrule
      Housing stock 2011--12 & 11plus12 & 7511 & 2011--12, 2012--13 & \checkmark & \checkmark & \checkmark \\
      Housing stock 2012--13 & 12and13  & 7802 & 2012--13, 2013--14 & \checkmark & \checkmark & \checkmark \\
      Fieldwork 2013--14     & fs13     & 7801 & 2013--14           & \checkmark &            &            \\
      Housing stock 2013--14 & 13plus14 & 8010 & 2013--14, 2014--15 &            &            & \checkmark \\
      Fieldwork 2014--15     & fs14     & 8009 & 2014--15           & \checkmark &            &            \\
      Housing stock 2014--15 & 14plus15 & 8186 & 2014--15, 2015--16 & \checkmark & \checkmark & \checkmark \\
      Housing stock 2015--16 & 15plus16 & 8350 & 2015--16, 2016--17 &            &            & \checkmark \\
      Housing stock 2016--17 & 16\_17   & 8494 & 2016--17, 2017--18 & \checkmark & \checkmark & \checkmark \\
      Housing stock 2019--20 & 19plus20 & 9058 & 2019--20, 2020--21 & \checkmark & \checkmark & \checkmark \\
      \bottomrule
    \end{tabular}
    \begin{tablenotes}[flushleft]\footnotesize
      \item \textit{Notes:} ``Homeownership'' and ``A--C share'' are the
      estimation targets (Table~\ref{tab:target-samples}); ``Regression''
      is the age-gradient regression of Figure~\ref{fig:ehs-age-gradient}.
      Fieldwork years run from April to March.
    \end{tablenotes}
  \end{threeparttable}
\end{table}

\paragraph{Definitions.} Age is the age of the household reference person
(HRP). All releases report it in bands (25--34, 35--44, 45--54, 55--64, and
65 and over); exact age is available only in 11plus12, 12and13, and fs13,
so statistics for ages 65--75 (inclusive) use those releases. A household is
an owner if it owns outright or with a mortgage. A dwelling is built by 1990
if the EHS dwelling-age variable places it in one of the bands before 1919,
1919--44, 1945--64, 1965--80, or 1981--90. A dwelling is in A--C if its
EPC band is A, B, or C; the EHS reports A and B as one band. The band is
based on SAP 2009 in 11plus12 and on SAP 2012 in the later releases; the
earlier release has only partial SAP 2012 recalculations, so we use its
published SAP 2009 band. Years since purchase are the completed years the
household has owned the home: exact in 11plus12 and 12and13, and in bands
(less than one, one, two, 3--4, 5--9, 10--19, 20--29, and 30 or more years)
in later releases. A recent buyer has owned the home for at most four
completed years.

\paragraph{Weights and pooling.} All EHS statistics use the household
weight supplied with each release. When a statistic pools releases, we
first normalize the weights so that they sum to one within each release and
then average across releases. Each release therefore receives equal weight,
whatever its sample size.

\paragraph{Overlap between releases.} Because each housing-stock release
pools two fieldwork years, adjacent releases contain the same households.
In our files, 11plus12 and 12and13 share 6,058 households (the 2012--13
fieldwork year). The fieldwork release fs13 repeats 2013--14 households
that also appear in 12and13, and fs14 repeats 2014--15 households that also
appear in 14plus15. Pooled point estimates count these households twice.
Our bootstrap standard errors (Section~\ref{sec:estimation}) draw each of
the 6,058 households shared by 11plus12 and 12and13 once. Household
identifiers in releases after 2013 are anonymized separately in each file,
so the other overlaps cannot be linked, and the standard errors of the
homeownership targets are somewhat understated.

% ----------------------------------------------------------------------
\subsection{The age gradient in the SAP rating.}\label{app:data-regression}

Figure~\ref{fig:ehs-age-gradient} reports estimates of
\begin{equation}\label{eq:ehs-age-gradient}
  \log \mathit{SAP}_i = \alpha + \sum_{j=2}^{5}\beta_j\,
  \ind\{\text{age group}_i = j\} + \gamma' X_i + \delta_{t(i)} + \eta_{r(i)}
  + \varepsilon_i,
\end{equation}
where $\mathit{SAP}_i$ is the SAP rating of household $i$'s dwelling, our
measure of energy efficiency; a higher rating means a more efficient
dwelling. The age groups are
35--44, 45--54, 55--64, and 65 and over, with 25--34 as the reference group.
$X_i$ contains log floor area, log equivalized household income before
housing costs (deflated by CPIH, 2015 prices), and indicators for a
single-person household, dwelling age, household size (top-coded at five),
and the HRP's labor-force status. $\delta_t$ and $\eta_r$ are release and
region (nine regions) fixed effects. The three specifications in the figure
are: age indicators only; adding $X_i$; and adding $X_i$ and the fixed
effects. The figure plots $\beta_j$; the text reports the percentage
difference $\exp(\beta_j)-1$. Standard errors are clustered
by dwelling age and region (81 clusters), and the intervals use the $t$
distribution.

The sample stacks the seven housing-stock releases marked ``Regression'' in
Table~\ref{tab:ehs-releases}, keeps owner-occupiers with an HRP aged 25 or
over, and contains 34,767 observations. A version that adds private renters
has 48,705 observations. The SAP rating is based on SAP 2009 in 11plus12 and
on SAP 2012 in later releases. The regressions are unweighted. For owners
aged 65 and over, the SAP rating is 9.0 percent lower than at 25--34 without
controls, 6.5 percent lower with $X_i$, and 6.6 percent lower with $X_i$ and
the fixed effects.

% ----------------------------------------------------------------------
\subsection{EPC register.}\label{app:data-epc}

\paragraph{Data.} We use the bulk download of all domestic certificates in
England and Wales from the Energy Performance of Buildings register, with one
file per local authority (343 authorities). Lodgement dates run from 1
October 2008 to 31 January 2023. The register contains 24,133,533
certificates and no duplicate certificate numbers.

Each certificate reports the dwelling's unique property reference number
(UPRN), the lodgement date, the current EPC band and SAP score, the
construction-age band, property type, floor area, modelled energy use, the
reason for the assessment (for example, a marketed sale or a new letting),
and the tenure stated by the assessor. A separate file lists recommended
improvements with indicative costs. We classify tenure from its text as
owner-occupied (12,153,496 certificates), rented (9,255,763), or
other or unknown (2,724,274). Construction-age bands are mapped to their
final year. Homes are built by 1990 if that year is 1990 or earlier, which
includes the bands ``before 1900'' and ``1983--1990''.

\paragraph{Consecutive pairs and annual rates (Table~\ref{tab:epc-transitions}).}
We keep owner-occupied certificates with a UPRN (12,006,376 certificates).
Of these, 425,058 have a construction-age band with no year (mostly
recorded as missing or invalid) and so have no vintage. We
order them by lodgement year within each UPRN, and pair each certificate
with the preceding one. Certificates of other tenures are removed before
pairing, so a pair can span an intervening letting. We keep pairs whose
lodgement years differ by at least one, leaving 1,705,904 pairs; a pair
lodged within a single calendar year is dropped and not linked further
back. A pair takes the vintage of its later certificate. For a pair $k$ with
gap $g_k\geq 1$ in calendar years, the annual upgrade rate is the mean over
pairs whose earlier certificate is in D--G of $\ind\{\text{later
certificate in A--C}\}/g_k$. The annual decay rate is the mean over pairs
whose earlier certificate is in A--C of $\ind\{\text{later certificate in
D--G}\}/g_k$. Dividing by the gap treats the probability of a change over
$g$ years as $g$ times an annual probability. This holds approximately when
changes are rare, and it misses changes that are reversed between two
certificates. Among homes built by 1990, the upgrade rate is 6.9 percent
(1,295,552 pairs) and the decay rate 15.2 percent (171,041 pairs). The
mean gap is 5.2 years. For homes built after 1990 the rates are 11.9
percent (104,678) and 5.8 percent (123,400).

\paragraph{Pairs around a sale.} The split by sale timing in
Section~\ref{sec:data-transitions} classifies pairs by the reason for
assessment recorded on each certificate, not by a match to transaction
records. A pair is \emph{at sale} if its later certificate was lodged for a
marketed or non-marketed sale. It is \emph{after a purchase} if its earlier
certificate was lodged for a sale, its later certificate was not, and the
two are at most two calendar years apart. All other pairs are \emph{away
from a sale}. This split is computed by a separate script that reads the
raw register and orders certificates within a year by their lodgement
time. Its pooled rates are 7.3 percent for upgrades and 11.2 percent for
decay, against 7.4 and 11.1 percent in Table~\ref{tab:epc-transitions}. For
homes built by 1990, the upgrade rate is 7.3 percent at sale (835,104
pairs), 6.6 percent after a purchase (133,087), and 5.5 percent away from a
sale (326,260). The decay rate is 10.7 percent at sale (124,902), 49.0
percent after a purchase (10,852), and 20.9 percent away from a sale
(36,616). Pairs after a purchase are short (1.4 years on average) and
few, and a sale certificate followed by a different rating within two
years points to measurement error at either assessment as much as to
physical change.

\paragraph{First and latest certificates.}
Figures~\ref{fig:epc-first-last} and~\ref{fig:epc-transition-matrix}
compare, for each owner-occupied dwelling with at least two certificates
at least one calendar year apart (1,557,929 dwellings), the first and the
latest certificate. The mean gap is 5.9 years. Among homes built by 1990
that start in D--G (1,178,403), 22.9 percent are in A--C at their latest
certificate. The share is 29.1 percent for homes starting in D and between
15 and 18 percent for homes starting in E, F, or G. Of the moves from D--G
into A--C, 95 percent end in band C.

\paragraph{Selection.} The register is not a random sample of the housing
stock. A certificate is lodged when a home is sold, let, or built, or when
an energy-efficiency scheme requires an assessment, so homes that change
hands often are over-represented, and the timing of a certificate is
itself a choice. Assessment conventions and default inputs change over
time and differ across assessors, and they can move a home across the C/D
boundary without physical work. We therefore use the register to describe
the direction and rough size of changes in efficiency status. The model is
not estimated on it.


\begin{figure}[htbp]
  \centering
  \includegraphics{fig-epc-first-last.pdf}
  \caption{Efficiency at the latest certificate by band at the first
  certificate, owner-occupied homes}
  \label{fig:epc-first-last}
  \begin{minipage}{0.96\textwidth}\footnotesize
    \textit{Notes:} For each EPC band at a dwelling's first owner-occupied
    certificate, the share of dwellings in bands A--C at the latest
    owner-occupied certificate. Dwellings with at least two such
    certificates at least one calendar year apart, 1,557,929 dwellings
    (mean gap 5.9 years). Left: all homes. Right: homes built by 1990,
    based on the latest certificate. EPC register, England and Wales,
    October 2008 to January 2023.
  \end{minipage}
\end{figure}

\begin{figure}[htbp]
  \centering
  \includegraphics{fig-epc-transition-matrix.pdf}
  \caption{Transitions between EPC bands, first to latest certificate,
  owner-occupied homes}
  \label{fig:epc-transition-matrix}
  \begin{minipage}{0.96\textwidth}\footnotesize
    \textit{Notes:} Row percentages: the distribution of the band at the
    latest certificate for each band at the first certificate. Lines mark
    the boundary between bands C and D. Sample as in
    Figure~\ref{fig:epc-first-last}.
  \end{minipage}
\end{figure}


% ----------------------------------------------------------------------
\subsection{Wealth and Assets Survey.}\label{app:data-was}

We use the household files of WAS rounds 5 to 8, which cover Great Britain
between April 2014 and March 2022. The
sample keeps households with a positive cross-sectional weight and drops
shared owners. Age is the HRP's age band; WAS does not release exact age.
Pooled statistics give each round equal weight. Negative asset values are
set to missing.

\paragraph{Housing share of wealth.} Housing wealth is the value of the
main residence for owners and zero for others. The target is the weighted
mean, over households aged 35--64 with positive gross wealth, of housing
wealth divided by the sum of housing wealth and gross non-housing financial
wealth excluding endowments. Both numerator and denominator are gross of
mortgage debt. Renters with financial wealth enter with a share of zero.

\paragraph{Liquid wealth to income.} Liquid wealth is gross non-housing
financial wealth excluding endowments, the closest harmonized household
measure in WAS. It includes all financial assets, not only cash and
deposits, and is not net of financial liabilities. Income is net regular
household income before housing costs: \texttt{DVTotNIRR5} in round 5 and
\texttt{DVTotInc\_BHC} in rounds 6--8. The targets are weighted medians of
the household-level ratio for HRPs aged 65--74 and 75 and over, among
households with positive income. The model counterparts are ages 65--75 and
75 to the terminal age. Medians of household ratios differ from ratios of
medians; the corresponding median incomes are £24,900 at 65--74 and
£18,200 at 75 and over, and median liquid wealth is £29,936 and £23,499.

\paragraph{Consumption and retirement income.} The consumption target uses
published expenditure rather than microdata. Its numerator is mean annual
household expenditure for an HRP aged 50--64 from ONS Family Spending
(Table A11, United Kingdom, financial years ending 2016 and 2017). We count
COICOP divisions 1--12 and spending on holidays abroad, but exclude
housing (division 4, apart from water and other dwelling services) and
moving costs; the model treats rent, energy, and maintenance separately. The weekly figures are multiplied by 52 and expressed in 2016
pounds with the CPI. The denominator is the weighted mean net household
income in WAS rounds 5 and 6 of households interviewed in the same two
financial years, with an HRP aged 50--64 (2,628 and 2,721 households). The
target is the ratio of the two equal-year means, 0.639. The
retirement-income target is the ratio of weighted mean net household income
at HRP ages 65--74 to that at 55--64 in the same WAS sample, 0.843; it
compares households, not a pension replacement rate for individuals.
Both targets use exact five-year HRP age bands from the person file. The
numerator covers the United Kingdom and the denominator Great Britain.


\paragraph{Equivalence scale.} The household-equivalence scale $e_j$ in the
model is the square root of mean household size by HRP age in WAS rounds
5--8. Means by age band are fitted with a cubic in age over ages 23--82,
and the scale is rescaled so that its mean over ages 35--64 matches a
reference value. 
% ----------------------------------------------------------------------
\subsection{Understanding Society.}\label{app:data-ukhls}

\paragraph{Income process.} The income process is estimated from waves 1--10
(2009--2019) of Understanding Society, restricted to England. Interviews from
wave 10 conducted in 2020 are excluded. The sample consists of male
household reference persons aged 20--64 who work 10--112 hours a week and
are not retired, students, caring for family, or long-term sick. It further
requires positive gross household income, hourly household earnings
between £2 and £400 in 2015 prices, and at least five observations per
person. This leaves 32,698 person-years. Income is gross monthly household
income times twelve, deflated to 2015 prices. The regressions are
unweighted. We regress log income on age indicators, marital status, the
number of children, and education with person fixed effects (31,711
observations, 4,265 persons). The deterministic age profile is a cubic in
age fitted to the predicted profile. The residual is the sum of a
permanent random walk and a transitory shock. Their variances come from
the growth of the residual's variance with the lag $d=1,\dots,9$: the
slope identifies the permanent variance (0.0092) and half the intercept
the transitory variance (0.0424). The initial permanent variance is the
cross-sectional variance at age 23 net of the two (0.168).

\paragraph{Owner mobility.} The annual moving rates of owners used as an
external check pair consecutive waves from 2--3 to 12--13. The base sample
is household reference persons in England who own outright or with a
mortgage at wave $t$, aged 23--90, with a positive cross-sectional weight.
A move means the person's address changed between $t$ and $t+1$, among
those observed at $t+1$. Weights are normalized to mean one within each pair
of waves. The rates are 9.6 percent at ages 25--34 (6,975 person-years),
3.6 percent at 35--64 (60,170), and 2.2 percent at 65--75 (22,605).

% ----------------------------------------------------------------------
\subsection{Mortality and energy bills.}\label{app:data-other}

\paragraph{Mortality.} Survival probabilities come from the ONS National Life
Tables for the United Kingdom, 2015--17, averaging the male and female
number of survivors $l_x$ at each age. The model uses the implied annual
survival probabilities for ages 23--81. Beyond the table, hazards are
extrapolated with a constant growth rate equal to the median of the last
seven hazard ratios, bounded between 1 and 1.25. The survival probability
in the recurring terminal state (82 and over) is one minus the inverse of
the expected number of remaining periods at 82, which gives 0.880.

\paragraph{Energy bills.} Energy bills by efficiency state are based on
owner-occupied certificates in the EPC register for homes built by 1990,
lodged between 2015 and 2019, using the latest certificate of each building
(2,833,945 buildings). We trim modelled consumption to 1--1,500 kWh per
square meter and floor area to 10--1,000 square meters. The register
reports modelled primary energy use. Unweighted means are 190.9 kWh per
square meter a year in bands A--C and 322.3 in D--G. We convert primary
to delivered energy by dividing by 1.59, the SAP 2012 primary-energy factors
(1.22 for gas, 3.07 for electricity) averaged with an assumed 80/20
gas/electricity mix. We then multiply by a common floor area of 101.75
square meters and a price of £0.0602 per kWh. The price is the 80/20 average
of DESNZ domestic unit prices, including VAT and excluding standing charges,
averaged over 2015--2019: £0.0381 for gas (Great Britain) and £0.1485 for
electricity (United Kingdom). Annual bills are £735 in A--C and £1,242 in
D--G.

% ----------------------------------------------------------------------
\subsection{Samples behind the estimation targets.}\label{app:data-targets}

Table~\ref{tab:target-samples} lists the source, sample, statistic, number
of observations, and number of releases or rounds behind each of the eleven
targets. Different targets describe different populations. The
homeownership rates are shares of all households. The A--C shares condition
on ownership and on homes built by 1990. The wealth ratios cover all
households in the relevant age band.

\begin{table}[htbp]
  \centering
  \caption{Samples behind the estimation targets}\label{tab:target-samples}
  \begin{threeparttable}
    \small
    \begin{tabularx}{\textwidth}{@{}>{\raggedright\arraybackslash}p{0.31\textwidth}
      l>{\raggedright\arraybackslash}X r r@{}}
      \toprule
      Target & Source & Sample and statistic & Obs. & Waves \\
      \midrule
      \multicolumn{5}{@{}l}{\textit{Homeownership}} \\
      \quad Ages 25--34 & EHS & All households; weighted owner share & 12,533 & 7 \\
      \quad Ages 35--64 & EHS & All households; weighted owner share & 46,055 & 7 \\
      \quad Ages 65--75 & EHS & All households, exact age; weighted owner share & 5,801 & 3 \\
      \addlinespace
      \multicolumn{5}{@{}l}{\textit{Owners in an A--C home, built by 1990}} \\
      \quad Ages 25--34 & EHS & Owners of homes built by 1990; weighted A--C share & 1,646 & 5 \\
      \quad Ages 35--64 & EHS & Owners of homes built by 1990; weighted A--C share & 12,279 & 5 \\
      \quad Ages 65--75 & EHS & As above, exact age & 1,831 & 2 \\
      \addlinespace
      \multicolumn{5}{@{}l}{\textit{Wealth, consumption and income}} \\
      \quad Housing share of wealth, 35--64 & WAS & Households with positive
        gross wealth; weighted mean of the household's housing share & 34,152 & 4 \\
      \quad Liquid wealth to income, 65--75 & WAS & HRP aged 65--74, positive
        income; weighted median ratio & 15,587 & 4 \\
      \quad Liquid wealth to income, 75+ & WAS & HRP aged 75 and over,
        positive income; weighted median ratio & 12,911 & 4 \\
      \quad Consumption to income, 50--64 & ONS, WAS & Mean core expenditure
        (published tables) over WAS mean net income & 5,349 & 2 years \\
      \quad Retirement income, 65--74 to 55--64 & WAS & Ratio of weighted mean
        net household income & 7,652 & 2 years \\
      \bottomrule
    \end{tabularx}
    \begin{tablenotes}[flushleft]\footnotesize
      \item \textit{Notes:} Observations are pooled households before
      equal-wave averaging. EHS waves are the releases in
      Table~\ref{tab:ehs-releases}; WAS waves are rounds 5--8. The last two
      targets use WAS households interviewed in the financial years ending
      in 2016 and 2017. The count for the consumption ratio refers to the
      income denominator.
    \end{tablenotes}
  \end{threeparttable}
\end{table}

\paragraph{Homeownership.} The three homeownership targets pool the seven
releases marked in Table~\ref{tab:ehs-releases} and include every household
with a positive weight, whatever its tenure. Ages 25--34 and 35--64 use the
banded age. Ages 65--75 use exact age and so rest on 11plus12, 12and13, and
fs13.

\paragraph{A--C shares.} The A--C targets use the five housing-stock releases
with a physical survey. The sample consists of owners of homes built by 1990
with a recorded EPC band. Ages 65--75 use exact age and so rest on 11plus12
and 12and13. These two releases cover fieldwork from 2011--12 to
2013--14, and the 2012--13 households appear in both. The recent-buyer shares
quoted in Section~\ref{sec:data} use the same five releases and the
same definitions. They are not targets.


% ======================================================================
